\documentclass[quick,con]{debate}
\motion{短视频提升了 / 降低了人的认知能力}
\stance{反方}{短视频降低了人的认知能力}
\begin{document}
\debatetitle
\begin{thesisbox}[一句话立论]短视频用每天两个半小时的日常，换走了人赖以深度思考的连续注意，这笔账是净亏的。\end{thesisbox}
\section{定义与判准（标准表述，全队同一口径）}
\begin{fields}
\field{短视频}{以手机为主要终端、单条时长通常在数分钟以内、由算法以信息流方式连续分发的视频内容，涵盖娱乐、生活、知识科普。（与正方共用中立定义，不要自创偏向定义）
}
\field{认知能力}{知觉、注意、记忆、理解、推理、判断与语言这一整套心理加工能力（美国心理学会《心理学词典》）。\textbf{补一句：这不是并列的七块砖，是一座楼，注意与执行控制在最下面。}
}
\field{判准}{在这一整套维度上比较普通人在短视频普及前后的整体状态，看\textbf{净方向}。\textbf{我方要证明}损耗压过增益；\textbf{对方要证明}增益压过损耗。
}
\field{对方提出偏向定义或判准时的 30 秒口径}{对方若把认知能力扩成“人加工具的总效能”：“辩题主语是人。给我一台计算器，我的数学能力不会提高。”对方若指控我方是“拔网线判准”：“我方不反对工具，只问它有没有让你自己加工。”
}
\end{fields}
\section{我方论点}
\begin{longtable}{@{}L{\dimexpr0.045\linewidth-1.50\tabcolsep}L{\dimexpr0.154\linewidth-1.50\tabcolsep}L{\dimexpr0.401\linewidth-1.50\tabcolsep}L{\dimexpr0.401\linewidth-1.50\tabcolsep}@{}}
\toprule \thead{标签} & \thead{一句话} & \thead{核心数据 / 学理 / 案例} & \thead{出处}\\\midrule\endhead
\textbf{注意力损耗} & 十几秒一条、无限下滑、划走就换——大脑学到的是切换有回报、忍耐没有回报 & 数据：2024 年一项对 48 名中国青年的问卷 + 注意网络测验 + 脑电研究，短视频成瘾倾向越高，冲突条件下额区脑电信号越弱（r = −0.395）。学理：注意网络模型（警觉、定向、执行控制），注意力是被环境训练出来的。案例：微短剧用户 6.62 亿人、占网民 59.7\%，“两三分钟一个反转”已成产业 & Yan et al. 2024, Frontiers in Human Neuroscience（已核实，\textbf{相关非因果，必须自己说}）；Posner \& Petersen 1990（有把握）；CNNIC 第 55 次报告（已核实）\\
\textbf{时段碎片化} & 它不只占掉时间，更把整块时间切碎，而深度理解只在连续时段里发生 & 数据：短视频应用人均单日使用时长 156 分钟（约 2.6 小时）× 短视频用户 10.40 亿人、占网民 93.8\%。学理：刻意练习——能力靠长时间、有难度、有反馈的持续投入，而练习需要不被打断的时间 & 《中国网络视听发展研究报告（2025）》，中国网络视听协会 2025 年 3 月发布（有把握：仅见媒体转述）；CNNIC 第 55 次报告（已核实）；Ericsson et al. 1993（有把握）\\
\bottomrule\end{longtable}
\section{对方论点 → 一句话反驳}
\begin{longtable}{@{}L{\dimexpr0.188\linewidth-1.33\tabcolsep}L{\dimexpr0.316\linewidth-1.33\tabcolsep}L{\dimexpr0.496\linewidth-1.33\tabcolsep}@{}}
\toprule \thead{对方会说} & \thead{我方一句话回} & \thead{对方追问时}\\\midrule\endhead
知识普惠（10.40 亿人、95\% 说学到了） & 你证明的是“接触到了”，辩题问的是“能力变了没有” & “从零到一难道不是增量？”→“如果他真学到了，是。你今天的证据只能证明他刷到了。请举出一项能力测验。”\\
形式适配（Guo 2014、Mayer 分段原则） & 你的证据来自有课程、有测验、自己选课的慕课，不是划到什么算什么的信息流 & “你承认短形式有效吗？”→“在有教学设计的场景里承认。那叫微课，不叫短视频。”\\
认知卸载 / 省下资源做高阶加工 & 省下来的资源，你没证明真被用在高阶加工上 & “记笔记不也是外包？”→“笔记是你主动写下、主动组织的，那本身就是加工。划视频不是。”\\
\bottomrule\end{longtable}
\section{必问与必答}
\begin{points}{必问 （对方不答就点名，自由辩前两分钟内问完）}
\item \textbf{请举出一项能力测验，显示短视频使用者成绩更好。}（全场核心）
\item 看到，和学到，是不是一回事？
\item 你那份 2014 年的研究，对象是课程视频还是信息流？
\item \textbf{一个人过去一天里有几段不被打断的一小时，今天还剩几段？}（一辩稿已抛出，全场追问）
\end{points}
\begin{points}{必答 （15 秒内正面回）}
\item 你证明的是有损伤还是净降低 → “净降低。受损的是底座，对方的增益只在‘看到’这一层。”
\item 48 人能推到十亿人吗 → “脑电是旁证。主力是时间与产品结构：每天两个半小时，覆盖 10.40 亿人。”
\item 你那份研究是相关还是因果 → “相关。我方一辩立论时就主动说明了。”
\item 知识可得性提高了吗 → “提高了。可得性不是能力。”
\end{points}
\section{自由辩战场概要}
\begin{longtable}{@{}L{\dimexpr0.176\linewidth-1.50\tabcolsep}L{\dimexpr0.286\linewidth-1.50\tabcolsep}L{\dimexpr0.231\linewidth-1.50\tabcolsep}L{\dimexpr0.307\linewidth-1.50\tabcolsep}@{}}
\toprule \thead{战场} & \thead{开场问题} & \thead{目标} & \thead{收束语}\\\midrule\endhead
一 · 接触与能力的鸿沟 & 请举出一项能力测验，显示短视频使用者成绩更好 & 让评委记住对方的数据都停在“看到”这一层 & “看到不等于学到，这一句是对方自己承认的。”\\
二 · 一天还剩几段完整时间 & 过去一天有几段不被打断的一小时，今天还剩几段？ & 让每位评委想起自己的一天 & “能力是练出来的，而练习需要不被打断的时间。”\\
三 · 这个工具让你加工了吗 & 书要你自己解码文字，短视频要你做什么？ & 拆掉“你反对一切工具”的元攻击，立起价值层 & “我方只问一句：这个工具，有没有让你自己加工过。”\\
\bottomrule\end{longtable}
时间分配 90 / 90 / 60 秒，后段不开新战场。站位：反一定义判准，反二证据，反三盘问成果，反四收束。

\section{如果……就……}
\begin{contingency}
\ifthen{对方定义不同，或拿出没预料到的数据}{“辩题主语是人”一句带回；新数据先问机构与年份，再问“它测的是‘看到’还是‘能力’？”}
\ifthen{论点一被打穿，或对方回避必问}{收缩到论点二：“就算注意力这一层打平，整块时间被切碎对方全场没回应过。”回避就点名、重复、不换题。}
\ifthen{红线（全队）}{脑电研究必须同一口气说“相关不是因果”；禁止说“注意力 8 秒不如金鱼”、禁止正面引用 Ophir 2009、禁止说“短视频让人变蠢”；挪威队列数据只能用来说明“群体认知水平会因环境逆转”，且必须交代它早于短视频。}
\end{contingency}
\end{document}
